\documentclass[12pt,a4paper]{article}
\newcommand{\CadernoDisciplina}{Análise Exploratória de Dados}
\newcommand{\CadernoCorHex}{17365D}
% Preâmbulo compartilhado dos cadernos de exercícios UNIFACOL.
% Definir \CadernoDisciplina e, opcionalmente, \CadernoCorHex antes deste input.
\providecommand{\CadernoDisciplina}{Disciplina}
\providecommand{\CadernoCorHex}{17365D}
\usepackage[a4paper,margin=2.05cm,headheight=15pt]{geometry}
\usepackage[T1]{fontenc}
\usepackage[utf8]{inputenc}
\usepackage[brazil]{babel}
\usepackage{lmodern,microtype,amsmath,amssymb}
\usepackage{enumitem,booktabs,tabularx,array,longtable,adjustbox,multirow}
\usepackage{xcolor,hyperref,graphicx,fancyhdr}
\usepackage[most]{tcolorbox}
\usepackage{tikz}
\usetikzlibrary{arrows.meta,positioning,fit,calc,patterns,shapes.geometric}
\definecolor{disciplinecolor}{HTML}{\CadernoCorHex}
\definecolor{stimulusgrey}{HTML}{EEF2F5}
\definecolor{answerorange}{HTML}{FFF0DE}
\definecolor{answerframe}{HTML}{B85C00}
\definecolor{supportgrey}{HTML}{F7F7F7}
\hypersetup{colorlinks=true,linkcolor=disciplinecolor,urlcolor=disciplinecolor}
\newcolumntype{Y}{>{\raggedright\arraybackslash}X}
\newcolumntype{C}{>{\centering\arraybackslash}X}
\setlength{\parindent}{0pt}
\setlength{\parskip}{.42em}
\setlist[enumerate]{leftmargin=*,itemsep=.28em,topsep=.35em}
\pagestyle{fancy}
\fancyhf{}
\lhead{UNIFACOL --- \CadernoDisciplina}
\rhead{Caderno ENADE --- 2026.2}
\cfoot{\thepage}
\newcommand{\ItemHeader}[4]{%
  \par\medskip\noindent
  \colorbox{disciplinecolor}{\parbox{\dimexpr\textwidth-2\fboxsep\relax}{\color{white}\bfseries Questão #1\hfill #2}}%
  \par\smallskip\noindent\textbf{Competência:} #3\par
  \noindent\textbf{Suporte:} #4\par\smallskip
}
\newcommand{\TextoBase}{\noindent\colorbox{stimulusgrey}{\parbox{\dimexpr\textwidth-2\fboxsep\relax}{\textbf{TEXTO-BASE}}}\par\smallskip}
\newcommand{\Comando}{\par\smallskip\noindent\textbf{COMANDO.} }
\newcommand{\FonteDidatica}{\par\smallskip\noindent{\footnotesize\textit{Fonte: situação e dados elaborados para fins didáticos.}}}
\newcommand{\FonteExterna}[1]{\par\smallskip\noindent{\footnotesize\textit{Fonte: #1}}}
\newcommand{\EspacoDiscursiva}{\par\noindent\rule{\textwidth}{.2pt}\vspace{1.15cm}\par\noindent\rule{\textwidth}{.2pt}}
\newtcolorbox{AnswerBox}{enhanced,breakable,colback=answerorange,colframe=answerframe,boxrule=.7pt,arc=1mm,title=Resposta explicada,fonttitle=\bfseries,coltitle=white,before skip=7pt,after skip=7pt}
\newtcolorbox{DocumentBox}[1][Documento ou artefato]{enhanced,breakable,colback=supportgrey,colframe=black!55,boxrule=.6pt,arc=.7mm,title={#1},fonttitle=\bfseries,coltitle=white,before skip=7pt,after skip=7pt}
\newtcolorbox{MemoBox}{enhanced,breakable,colback=white,colframe=disciplinecolor,boxrule=.7pt,arc=.7mm,title=Memorando,fonttitle=\bfseries,coltitle=white,before skip=7pt,after skip=7pt}
\newcommand{\LegendaDidatica}[1]{\par\smallskip{\footnotesize\textit{#1. Fonte: elaborado para fins didáticos.}}}

\setcounter{secnumdepth}{0}
\begin{document}
\begin{center}
{\Large\bfseries UNIFACOL}\\[2mm]
{\large Avaliação Final}\\[1mm]
{\Large\bfseries VERSÃO B}
\end{center}
\begin{DocumentBox}[Identificação]
\begin{tabularx}{\textwidth}{@{}p{.48\textwidth}p{.48\textwidth}@{}}
Estudante: \hrulefill & Turma: \hrulefill\\[3mm]
Professor: Cayo Oliveira & Data: \rule{2.7cm}{.2pt}/2026
\end{tabularx}
\end{DocumentBox}
\textbf{Instruções.} Esta avaliação contém seis questões objetivas e duas discursivas. Nas objetivas, assinale uma única alternativa. Nas discursivas, mostre cálculos intermediários e justifique decisões com as evidências do texto-base. Respostas sem desenvolvimento podem não receber pontuação integral.
\bigskip
% origem: AED-C05-O07
\ItemHeader{1}{Objetiva}{Aplicar linhagem e impacto}{Mudança de regra}
\TextoBase O comitê propõe mudar ``cliente ativo'' de compra nos últimos 90 dias para compra nos últimos 60 dias. O indicador alimenta três painéis, uma campanha automática e um modelo. A alteração pode mudar contagens e decisões mesmo que nenhuma linha bruta seja modificada, exigindo análise de impacto antes da vigência.
\FonteDidatica
\Comando Antes de publicar, a equipe deve
\begin{enumerate}[label=\Alph*)]
  \item alterar a SQL e aguardar reclamações.
  \item recalcular apenas dados futuros e misturar definições na mesma série.
  \item manter o mesmo nome sem registrar a quebra semântica.
  \item usar linhagem para identificar consumidores, versionar regra e vigência, comparar impacto e comunicar migração.
  \item atualizar somente o painel executivo por ser mais visível.
\end{enumerate}

% origem: AED-C13-O05
\ItemHeader{2}{Objetiva}{Escolher memória apropriada}{Correção de preferência e regra}
\TextoBase
Um usuário prefere gráficos em barras; a política corporativa determina que valores inferiores a cinco sejam suprimidos. O protótipo grava ambas as informações como lembranças livres produzidas pelo modelo.
A preferência de formato pertence ao usuário, mas uma política revogada não pode persistir como memória nem ser reintroduzida pelo histórico.
\FonteDidatica
\Comando O redesenho adequado é
\begin{enumerate}[label=\Alph*)]
  \item armazenar preferência e política juntas em memória conversacional sem versão.
  \item permitir que o agente escolha qual regra lembrar conforme a tarefa.
  \item tratar preferência como dado de perfil revogável e política como regra governada, versionada e aplicada fora da memória livre.
  \item converter a política em mensagem do usuário para que possa ser substituída em cada conversa.
  \item descartar a preferência e manter somente o histórico integral indefinidamente.
\end{enumerate}

% origem: AED-C03-D15
\ItemHeader{3}{Discursiva}{Escolher distribuição}{Planejamento operacional}
\TextoBase Uma equipe modelará quatro fenômenos: número de falhas em lote fixo sob condições estáveis; sucesso de uma única tentativa; instante de chegada em janela projetada sem preferência; e erro de medição aproximadamente simétrico produzido por muitos efeitos pequenos. A escolha da distribuição deverá decorrer do mecanismo e indicar o parâmetro e a premissa que pode falhar.
\FonteDidatica
\Comando Associe Bernoulli, binomial, uniforme ou normal a cada caso, explicando parâmetros, condições e uma razão que poderia invalidar cada escolha.
\EspacoDiscursiva

% origem: AED-C12-O05
\ItemHeader{4}{Objetiva}{Gerenciar estado e reexecução}{Duplo envio no Streamlit}
\TextoBase
Ao alterar um filtro, o Streamlit reexecuta o script. O protótipo dispara novamente uma inferência cara e grava duas decisões idênticas. Não há chave que relacione pergunta, dados e versão do modelo.
O incidente ocorre quando a interface reexecuta o script; a correção deve impedir duplicidade sem perder o histórico visível ao usuário.
\FonteDidatica
\Comando A correção adequada é
\begin{enumerate}[label=\Alph*)]
  \item remover os logs duplicados sem impedir o efeito repetido.
  \item desativar todos os widgets após a primeira execução da sessão.
  \item usar estado de sessão, formulário para submissão explícita, chave idempotente e cache versionado pelos insumos relevantes.
  \item guardar apenas a última resposta em variável global compartilhada por usuários.
  \item aumentar o timeout para que as duas chamadas terminem na mesma ordem.
\end{enumerate}

% origem: AED-C08-O05
\ItemHeader{5}{Objetiva}{visual}{Situação profissional e suporte quantitativo}

\TextoBase

Um mapa de calor compara bairros, mas não apresenta escala, usa cores sem ordenação e mistura contagem de atrasos com percentual de entregas. Bairros possuem volumes muito diferentes. A revisão deve preservar unidade, denominador e acessibilidade antes de destacar uma região.

\FonteDidatica

\Comando Selecione o redesenho que separa medidas, explicita escala e permanece legível sem depender apenas de cor.

\begin{enumerate}[label=\Alph*)]
  \item retirar legenda
  \item aumentar o número de cores
  \item ordenar bairros alfabeticamente apenas
  \item separar medidas, explicitar escala/unidade e oferecer rótulos acessíveis
  \item somar contagem e percentual
\end{enumerate}

% origem: AED-C14-D15
\ItemHeader{6}{Discursiva}{Responder a contestação e corrigir o sistema}{Decisão baseada em número errado}
\TextoBase
Um relatório GenBI levou a reduzir estoque. Depois, descobriu-se que o snapshot tinha atraso de dois dias e o trace não registrava versão. Houve perda operacional, mas a extensão ainda está em apuração.
A recomendação equivocada já afetou clientes; correção técnica deve ocorrer junto com contenção, comunicação, reparação e prevenção de recorrência.
\FonteDidatica
\Comando Proponha contenção, correção da decisão, preservação de evidências, comunicação de fatos e hipóteses, reconstrução, compensação/encaminhamento, correção de versionamento e gates para retorno.
\EspacoDiscursiva

% origem: AED-C02-O05
\ItemHeader{7}{Objetiva}{Aplicar média ponderada}{Notas com pesos}
\TextoBase O plano de avaliação atribui peso 2 a uma atividade prática e peso 8 à prova da unidade. Um estudante obteve, respectivamente, notas 9 e 6. O sistema precisa consolidar a nota respeitando a contribuição definida para cada componente, e não tratar as duas observações como igualmente representativas.
\FonteDidatica
\Comando A média que respeita os pesos é
\begin{enumerate}[label=\Alph*)]
  \item $(9+6)/10=1{,}5$.
  \item $(9+6)/2=7{,}5$.
  \item $(2\times9+8\times6)/10=6{,}6$.
  \item $(8\times9+2\times6)/10=8{,}4$.
  \item $(2+8)/(9+6)=0{,}67$.
\end{enumerate}

% origem: AED-C09-O09
\ItemHeader{8}{Objetiva}{alucinação}{Situação profissional e suporte quantitativo}

\TextoBase

A resposta cita uma tabela inexistente, embora a redação seja coerente e os números pareçam plausíveis. Não há artefato que sustente a afirmação. Antes de publicar, a equipe precisa exigir fonte verificável ou recusar a conclusão.

\FonteDidatica

\Comando Selecione a resposta operacional compatível com uma citação sem fonte existente.

\begin{enumerate}[label=\Alph*)]
  \item aumentar criatividade
  \item exigir referência válida e sustentação de cada alegação; caso contrário, abster-se
  \item treinar imediatamente com a resposta
  \item ocultar a citação
  \item aceitar pela coerência
\end{enumerate}

\end{document}
